\section*{الفصل \ref{ch:b3:learning-memory} --- اللدونة العصبية، والتعلم، والذاكرة}

\begin{solution}{exo:b3:learning-memory:1}
أُلغيت \omterm{def:b3:learning-memory:kinds}{الذاكرة التصريحية} (الوقائع والأحداث) عند H.M. --- إذ
لم يعد يقدر على استذكار شيء جديد --- بينما نجت الذاكرة غير
التصريحية: فقد تعلم الرسم في المرآة تعلمًا سويًّا من غير أن
يذكر الجلسات. والمهارات: المخطط \omterm{def:b3:nervous-systems:plan}{والمخيخ}؛ والخوف المشروط:
اللوزة؛ والمنعكسات الحركية المشروطة: \omterm{def:b3:nervous-systems:plan}{المخيخ}؛ والتهيئة:
القشرة الحسية.
\end{solution}

\begin{solution}{exo:b3:learning-memory:2}
حين يسهم عصبون قبل مشبكي مرارًا في إشعال عصبون بعد مشبكي
تتقوى وصلتهما. ومن ثم تكون \omterm{def:b3:learning-memory:ltp}{التقوية طويلة الأمد} نوعية الدخل
(فلا يتغير إلا المشابك الفعالة، التي شاركت)، وترابطية (فدخل
ضعيف فعال بينما تشتعل الخلية تحت دخل قوي يتقوى)، \omterm{def:b3:structural-biology:allostery}{وتعاونية}
(إذ يجب أن تعمل دخلات كافية معًا لإشعال الخلية).
\end{solution}

\begin{solution}{exo:b3:learning-memory:3}
فهو لا ينقل إلا حين يكون الغلوتامات مرتبطًا (أي إن الخلية
قبل المشبكية اشتعلت) \emph{و}يكون الغشاء زائل الاستقطاب بما
يكفي لطرد المغنيزيوم (أي إن الخلية بعد المشبكية فعالة):
فشرطا \omterm{def:b3:learning-memory:ltp}{قاعدة هِب} في بروتين واحد. وAP5 قبل الكزاز: فلا دخول
كالسيوم ولا تقوية، مع أن النقل عبر مستقبلات AMPA سويّ. وAP5
بعده: فالتقوية المستحدَثة فعلًا لا تتأثر --- فالمستقبل لازم
للاستحداث لا للصيانة.
\end{solution}

\begin{solution}{exo:b3:learning-memory:4}
القصير الأمد: تعديل تساهمي لبروتينات قائمة --- إذ يفسفر PKA
قناة بوتاسيوم ويغلقها، فتتسع الشوكة ويزداد الإطلاق؛ ولا
بروتين جديد. والطويل الأمد: يفعّل PKA عاملَ CREB، فتُستنسخ
مورّثات جديدة وتنمو نهايات مشبكية جديدة. ومثبط تركيب
البروتين المعطى أثناء التدريب المتكرر يترك التسهيل القصير
الأمد سليمًا ويلغي الذكرى المقيسة بعد أيام؛ وإذا أُعطي
بعده لم يفعل شيئًا.
\end{solution}

\begin{solution}{exo:b3:learning-memory:5}
$0.7^{n} \le 0.1$: أي $n \ge 6.5$، فتلزم $7$ محاولات. ومن $0.9\lambda$،
$0.9\times 0.7^{n} < 0.1$: أي $7$ محاولات أيضًا. والتناظر لأن
$\alpha$ نفسه يحكم الاتجاهين. وفي الحيوانات يكون الانطفاء
أبطأ عادةً وتعود الذكرى الأصلية بعد راحة (التعافي العفوي):
فالانطفاء تعلم جديد يكبت القديم، لا محو له.
\end{solution}

\begin{solution}{exo:b3:learning-memory:6}
المتجهان الذاتيان $(1,1)/\sqrt{2}$ بقيمة ذاتية $1.5$
و$(1,-1)/\sqrt{2}$ بقيمة $0.5$؛ وتدير \omterm{def:b3:learning-memory:ltp}{قاعدة هِب} المتجه
$\mathbf{w}$ نحو $(1,1)$، فيصير العصبون كاشفًا لاشتعال
الدخلين معًا. وخطوة أويا: $y = 1$،
$\Delta\mathbf{w} = 0.1\times 1\times((1,1) - 1\times(1,0)) = (0,0.1)$،
فيكون $\mathbf{w} = (1, 0.1)$.
\end{solution}

\begin{solution}{exo:b3:learning-memory:7}
$500\times 0.05 = 25$ أيونًا حرًّا؛ و$25/6\times 10^{23} = 4.2\times 10^{-23}$
مول في $10^{-16}$ لتر: أي \qty{0.42}{\micro mol/L}، أربعة أمثال
مستوى الراحة لكنها دون $K_{d}$ الكالمودولين --- أي تفعيل
جزئي فحسب؛ وتلزم عدة انفتاحات قنوية معًا لبلوغ المدى
الميكرومولاري الذي يشغّل CaMKII.
\end{solution}

\begin{solution}{exo:b3:learning-memory:8}
تُدرَّب A حتى $V_{A} = \lambda$: فتندفق الدوبامين في
المحاولات المبكرة وتصمت مع صيرورة المكافأة متنبَّأً بها.
وفي المحاولات المركّبة يكون التنبؤ $V_{A} + V_{B} = \lambda$
دقيقًا أصلًا، فالخطأ صفر، والدوبامين صامت، ولا تتغير $V_{B}$
أبدًا: فقد حُجبت B. وإذا اختُبرت وحدها لم تستثر استجابة؛
فإن تلت مكافأةٌ B وحدها كانت غير متوقعة فاندفق الدوبامين.
\end{solution}

\begin{solution}{exo:b3:learning-memory:9}
(أ) لا تقوية في \omterm{def:b3:learning-memory:hippocampus}{CA1} وعجز في الذاكرة المكانية، مع سلامة سائر
التعلم. (ب) الذاكرة القصيرة الأمد سوية، والطويلة الأمد
غائبة عند \qty{24}{h}. (ج) لا أثر --- \omterm{def:b3:learning-memory:kinds}{فالترسيخ} انتهى (إلا إذا
أُعيد تفعيل الذكرى، فتصير هشّة من جديد). (د) يتجمد الفأر في
السياق الآمن: فالذكرى تُستدعى اصطناعًا \omterm{def:b3:learning-memory:hippocampus}{بإعادة} تفعيل الخلايا
التي رمّزتها.
\end{solution}

\begin{solution}{exo:b3:learning-memory:10}
$\Delta|\mathbf{w}|^{2} \approx 2\mathbf{w}\cdot\Delta\mathbf{w} = 2\eta
y^{2} \ge 0$: فالطول لا ينقص أبدًا وينمو كلما اشتعل العصبون.
وعند نقطة ثابتة لأويا يكون $C\mathbf{w} = (\mathbf{w}^{\mathsf{T}}C
\mathbf{w})\mathbf{w}$، فيكون $\mathbf{w}$ متجهًا ذاتيًا عند $\lambda =
\lambda|\mathbf{w}|^{2}$، ومن ثم $|\mathbf{w}| = 1$. والنمو غير
المحدود يشبع كل مشبك، ويلغي الانتقائية، ويسوق تهييجًا
جامحًا؛ والمعيِّرات البيولوجية هي التحجيم المشبكي (إذ تُخفَّض
مشابك العصبون كلها حين يكون مفرط النشاط)، \omterm{def:b3:learning-memory:ltp}{والاكتئاب طويل
الأمد}، وخفض المشابك أثناء النوم.
\end{solution}

\begin{solution}{exo:b3:learning-memory:11}
الجرذ: $0.14\times 3\times 10^{5} \approx 4\times 10^{4}$؛ والإنسان: $0.14
\times 2\times 10^{6} \approx 3\times 10^{5}$. والتقدير يفترض
أنماطًا عشوائية كثيفة واتصالًا تامًا؛ والأنماط الحقيقية
متناثرة (وهو ما يرفع السعة كثيرًا) ومبنيّة، \omterm{def:b3:learning-memory:kinds}{والحُصين} مخزن
مؤقت يسلّم الذكريات إلى قشرة فيها $10^{10}$ عصبون
--- فمخزن العمر هو القشرة، ورقم \omterm{def:b3:learning-memory:kinds}{الحُصين} حجم مخزن مؤقت، لا
عدد ذكريات.
\end{solution}

\begin{solution}{exo:b3:learning-memory:12}
دخلات العين المفتوحة مترابطة مع اشتعال الخلية القشرية
فتتقوى؛ ودخلات العين المغلقة ليست كذلك، فتضعف تحت التنافس
على وزن كلي محدود --- أي هِب مع تعيير. وحين تُغلق العينان
كلتاهما لا يُحابى أي دخل، فيحتفظ كلاهما بإقليم (ضعيف) ويكون
الأثر أخف. وتُغلق الفترة حين تنضج الدارات المثبِّطة، وتتصلب
الشباك حول العصبونات عند المشابك، وتتبدل وحيدات \omterm{prop:b3:learning-memory:nmda}{مستقبل NMDA}
إلى صورة أقل تساهلًا.
\end{solution}

\begin{solution}{pb:b3:learning-memory:1}
\textbf{1.} $10\times 10^{3} = 10^{4}$ أيون؛ و\qty{2}{\%} حرة، أي $200$؛
و$200/6\times 10^{23} = 3.3\times 10^{-22}$ مول في $8\times 10^{-17}$ لتر:
أي \qty{4.2}{\micro mol/L}.
\textbf{2.} نعم، فوق عتبة \qty{2}{\micro mol/L}. وبمستقبلين:
$40$ أيونًا حرًّا، أي \qty{0.83}{\micro mol/L} --- دونها.
\textbf{3.} عند \qty{0.83}{\micro mol/L} يكون الكالسينورين فعالًا
وCaMKII غير فعال: فتُزال مستقبلات AMPA ويضعف المشبك (اكتئاب).
أيون واحد، وإنزيمان مختلفا الألفة: فالارتفاع الكبير السريع
يبلغ الكيناز، والصغير البطيء يبلغ الفسفاتاز وحده.
\textbf{4.} دخل ثانٍ في غضون بضعة ثوابت زمن يضيف كالسيومه إلى
ما بقي من الأول؛ وبعد \qty{50}{ms} لا يكاد يبقى شيء.
فاشتراط التزامن في NMDA وخمود \qty{20}{ms} يعطيان معًا نافذة
قدرها نحو \qty{20}{ms} في كل جهة --- أي نافذة توقيت الشوكات.
\textbf{5.} تلتقط الموقيات معظم الكالسيوم في جزء من ميكرومتر
ويبطئ عنق الشوكة الضيق الانتشار إلى التغصن، فيبقى الارتفاع
في الشوكة التي كانت فعالة؛ ولا ترى جارة على بعد
\qty{0.5}{\micro m} شيئًا، فلا يتغير إلا المشبك الفعال.
\textbf{6.} $(70 - 30)/5 = 8$ دخلات متزامنة، أو كمون عمل راجع
من جسم الخلية: فالمشبك الواحد لا يقدر على فك حجب مستقبلات
NMDA الخاصة به --- والتعاون مبنيّ في الفيزياء.
\textbf{7.} $(1,1)/\sqrt{2}$ بقيمة ذاتية $1.6$؛ و$(1,-1)/\sqrt{2}$
بقيمة $0.4$.
\textbf{8.} $C\mathbf{w}_{0} = (1, 0.6)$، و$\mathbf{w}_{1} = (1.10,
0.06)$، والزاوية إلى $(1,1)$: $45^{\circ} - 3.1^{\circ} = 41.9^{\circ}$.
و$\mathbf{w}_{2} = (1.21, 0.13)$: $38.8^{\circ}$. و$\mathbf{w}_{3} = (1.34,
0.22)$: $35.8^{\circ}$.
\textbf{9.} $w_{2}/w_{1} \to 1$ بينما ينمو $|\mathbf{w}|$ بلا حد
(بعامل $1.16$ في كل خطوة).
\textbf{10.} $C\mathbf{w} = (1.28, 1.28)$؛ و$\mathbf{w}^{\mathsf{T}}C
\mathbf{w} = 2.05$؛ و$\Delta\mathbf{w} = 0.1\times(1.28 - 2.05\times 0.8)
= -0.036$ لكل مركّبة: فيكون $\mathbf{w} = (0.76, 0.76)$، متقلصًا نحو
متجه الوحدة $(0.71, 0.71)$.
\textbf{11.} المشترك $y = 1.42$؛ والدخل 2 وحده $y = 0.71$، أي النصف.
\textbf{12.} الدخل الضعيف يشتعل بينما يسوق القوي الخلية، فيكون
مترابطًا مع الخرج وتقوّيه \omterm{def:b3:learning-memory:ltp}{قاعدة هِب} --- فيتعلم العصبون
التلازم.
\textbf{13.} $1 - 0.85^{n}$: أي $0.56$، و$0.80$، و$0.96$.
\textbf{14.} $0.85^{n} \le 0.05$: أي $n \ge 18.5$، أي $19$ محاولة.
\textbf{15.} لا: فنسبة المقارب المبلوغة بعد $n$ محاولة لا
تتوقف على $\lambda$؛ وإنما يتضاعف المقارب. وقد ترفع مكافأة
أكبر قيمة $\alpha$ (البروز)، وتُبلغ عتبة سلوكية على القيمة
\emph{المطلقة} أسرع.
\textbf{16.} $V_{B} = 0$: أي محجوبة. فالمركّب متنبَّأ به
تمامًا، والخطأ صفر، وعصبونات الدوبامين صامتة.
\textbf{17.} B وحدها لا تتنبأ بشيء، فالخطأ $\lambda$ ويندفق
الدوبامين؛ ثم ترتفع $V_{B}$ كما في الاكتساب، $\lambda(1 -
0.85^{n})$.
\textbf{18.} يوصل الدواء الدفقة التي تعني ``أفضل مما تُنبِّئ''
بعد أي شيء سبقها؛ فتزيد القاعدة قيمة تلك الإشارات والأفعال
في كل مرة، فتكتسب قيمة مستقلة عن أي نتيجة حقيقية ---
فيصير الطلب قهريًا.
\textbf{19.} $0.14\times 3\times 10^{5} \approx 4.2\times 10^{4}$.
\textbf{20.} $4.2\times 10^{4}/20 = 2100$ يوم، أي نحو ست سنين؛
فلا بد أن ينقل \omterm{def:b3:learning-memory:kinds}{الترسيخ} الذكريات إلى القشرة وأن يفرّغ النسيان
المخزن.
\textbf{21.} كتابة أنماط جديدة بسرعة في القشرة ستطمس المشابك
التي ترمّز القديمة (التداخل الكارثي)؛ وأما \omterm{def:b3:learning-memory:hippocampus}{الإعادة} البطيئة
المتشابكة فتتيح للقشرة أن تدمج الجديد بالقديم، بينما يحفظ
\omterm{def:b3:learning-memory:kinds}{الحُصين} الذكرى الجديدة في تلك الأثناء --- أي مخزن مؤقت سريع
يغذي مخزنًا بطيئًا.
\textbf{22.} $0.03\times 3\times 10^{5} = 9000$ خلية لكل مكان.
والأنماط المتناثرة تتداخل قليلًا، فيمكن تخزين عدد أكبر منها
بكثير قبل أن يفسد تداخلها الاسترجاع؛ وترتفع السعة أسرع من
$N$ كلما ازداد تناثر الأنماط، بثمن أن يستعمل كل نمط نوعية
خلايا أكثر.
\textbf{23.} الضغط بعامل $20$: فالخلايا التي تشتعل متباعدةً
\qty{400}{ms} في الركضة تشتعل متباعدةً \qty{20}{ms} في
\omterm{def:b3:learning-memory:hippocampus}{الإعادة} --- أي داخل نافذة توقيت الشوكات. ومغزى الضغط أن
يجلب متواليات أبطأ من أن تترابط في الزمن الحقيقي إلى متناول
\omterm{def:b3:learning-memory:ltp}{قاعدة هِب}.
\textbf{24.} بعد يوم تكون الذكرى لا تزال متوقفة على \omterm{def:b3:learning-memory:kinds}{الحُصين}؛
وبعد شهر تكون قد رُسّخت في القشرة ولم تعد تحتاج إليه ---
وهو بالضبط الفقد المتدرج زمنيًا عند H.M.، حيث تذهب الذكريات
الحديثة وتسلم البعيدة.
\textbf{25.} نحو \qty{4.2}{\micro mol/L} من الكالسيوم الحر بعد
عشرة انفتاحات؛ و$19$ محاولة لبلوغ \qty{95}{\%}؛ ونحو
$4\times 10^{4}$ نمط في \omterm{def:b3:learning-memory:hippocampus}{CA3} عند الجرذ.
\end{solution}
